\section{Limitations}\label{sec:disc:limitations}

All tabulated scores come from one synthetic corpus with full per-pulse labels.
The code is proprietary and is not released with this thesis.

The test modes are the training modes.
Held-out modulation types were not scored, even though clustering rather than a softmax is what would make that experiment possible.
Claims about novel emitters or novel modes at inference therefore go beyond the tables.
Emitter identifiers are recording-local by construction, which is the right operational assumption, but it is not a measured open-set test.

Every score is a window score.
The encoder never sees a full recording in one forward pass, trailing segments shorter than $L$ are dropped, and nothing is tracked across window boundaries.
Self-attention and the physical bias both scale as $L^{2}$. Extending the window would raise the memory gap in \cref{sec:disc:attention} and would still not be streaming inference.

The supervised contrastive kernel of \cref{eq:method:supcon} is a second quadratic term.
A batch of $B$ windows of length $L$ is pooled as $\mathcal{B}$ with $\lvert\mathcal{B}\rvert=BL$, and each branch forms all-pairs similarities $s_{ij}=\vect{z}_i^{\top}\vect{z}_j$ over that set, so the naive kernel is a $(BL)\times(BL)$ matrix per branch.
With $B=4$ and $L=2000$ that is $8000\times 8000$ per branch, and two branches.
The compute is $\mathcal{O}((BL)^{2})$ per term and cannot be avoided for this kernel, because every other pulse in the batch is a candidate in the denominator.
Memory need not hold the full matrix: the sum can be chunked.
A naive materialization does store $(BL)\times(BL)$.
This thesis does not report a measured loss-tensor size, unlike the Bias dump of \cref{sec:disc:attention}.
Centroid methods such as \gls{dec} and \gls{idec} score $\mathcal{O}(NK)$ against $K$ centroids, a softmax scores $\mathcal{O}(NC)$ classes, triplets sample a few pairs, and SwAV uses a small codebook.
Those losses avoid the all-pairs batch matrix.
Supervised contrastive needs the full candidate set to define positives by label and to prevent collapse.
Switching kernel would change the method, not just the implementation.

Pulse counts per label are unbalanced inside windows, as \cref{sec:exp:datasets} notes.
Rare modes and sparse emitters contribute little to a window-mean $\ARI$.
A method can look strong in \cref{tab:exp:attn_md} while still failing the labels that occupy few pulses.
The three-pulse emitter in \cref{fig:exp:feat_imbalance} and the crowded window in \cref{fig:exp:feat_maxmode} are reminders of that regime, not a substitute for a stratified breakdown.

The two terms were not reweighted.
The finding that joint training costs emitter $\ARI$ on Vanilla is therefore tied to the unweighted sum and to that backbone.
RoPE-TOA does not show the same cost: its joint emitter $\ARI$ already matches an emitter-only copy of that encoder (\cref{sec:exp:attention,tab:app:ari_em}).
Emitter-only Bias does: it falls to the RoPE-TOA+Bias band rather than staying at the joint Bias ceiling.

\Gls{dbscan} is part of the reported system.
The radius is chosen on the validation grid, $\mathrm{minPts}$ is fixed, and the noise category was ignored because it stayed small on this corpus.
A different density clusterer, or a different $\varepsilon$ rule, would change the discrete partitions without changing the embeddings.

Finally, no operational \gls{elint} recordings enter the numbers.
Simulator drop and spurious-pulse models exist in the generator, but they are not the test set behind the tables.
Distribution shift, library mismatch, and deinterleaving errors of the kind studied by \textcite{jakobsson2024trackassociation} and \textcite{coleman2023rwrdrift} remain outside the present scores.
